\documentclass[11pt]{article}

\usepackage{color,hyperref,amsmath}
\usepackage[sc,osf]{mathpazo} % Palatino font, w/ smallcaps & old-style figures
\usepackage[a4paper, total={5in, 8.25in}]{geometry}

\usepackage[utf8]{inputenc}
\usepackage{graphicx}
\usepackage{longtable}
\usepackage{enumitem}
\usepackage{tikz}
\usetikzlibrary{shadows.blur, positioning}

\newcommand{\code}[1]{\texttt{#1}}
\newcommand{\cmd}[1]{\texttt{#1}}

% The 5in measure is narrow for unbreakable \code identifiers; let TeX
% stretch interword space rather than overrun the margin.
\setlength{\emergencystretch}{1.5em}

% Default \tabcolsep (6pt each side of every column) adds up fast in
% this narrow 5in measure -- e.g. a 4-column table pays 48pt (0.67in)
% of padding alone. Trim it once, globally, rather than fighting
% column-width arithmetic per table.
\setlength{\tabcolsep}{4pt}

% Simple boxed CAUTION callout, styled plainly (no tcolorbox dependency).
\newcommand{\caution}[1]{%
  \par\noindent
  \fbox{\parbox{4.7in}{\textbf{CAUTION:} #1}}
  \par\smallskip
}

\title{WHAM-XREX-PFMG474\\ \large Firmware Build, Flash, and Verification \\ Standard Operating Procedure}
\author{}
\date{August 2026}

\begin{document}
\maketitle

% longtable, not tabular/center: this table grew past one page as of
% Rev E (five rows now), and a plain tabular cannot break across a page
% boundary -- same overflow bug already fixed once in the Troubleshooting
% table (Section~\ref{sec:troubleshooting}), below. \endhead repeats the
% header row on each continuation page.
\begin{longtable}{p{0.45in} p{0.75in} p{0.5in} p{2.7in}}
  \hline
  \textbf{Rev} & \textbf{Date} & \textbf{Author} & \textbf{Description} \\
  \hline
  \endhead
  A & 2026-08-31 & -- & Initial issue: build, bootstrap flash, serial
    reflash. No power-stage procedures yet (Section~\ref{sec:scope}). \\
  B & 2026-09-04 & -- & Added Section~\ref{sec:table-fire}: uploading a
    PFM table and firing it -- this firmware now drives real HRTIM PWM
    output (Section~\ref{sec:scope} updated accordingly). Serial baud
    raised 9600~$\to$~115200 throughout (Section~\ref{sec:verify} and
    elsewhere). \\
  C & 2026-09-22 & -- & Added Section~\ref{sec:profile-fire}: the modern
    closed-loop shot-profile procedure (\code{ARM}/\code{SHOT:*}),
    driven by hand over raw SCPI, bypassing \code{wham\_console.py}'s
    \code{shot} wizard. Corrected Section~\ref{sec:scope}'s Rev~B claim
    that this firmware has ``no ARM precondition, no state machine, no
    fault interlock'' -- all three now exist (added between Rev~B and
    Rev~C; Section~\ref{sec:table-fire}'s own caution still applies
    as-is to the legacy \code{TABLE:*}/\code{FIRE} path specifically,
    which remains interlock-free by design). \\
  D & 2026-09-22 & -- & Added Section~\ref{sec:telemetry}: reading and
    translating the \code{!EVT} telemetry stream and the retained
    flight-recorder log (\code{SYS:EVLOG?}). Updated
    Section~\ref{sec:profile-fire}: the external-enable
    (\code{EXTernal:ENAble}, PF13) and external-trigger
    (\code{EXTernal:TRIGger}, PF15) interlocks now default \textbf{ON}
    at every boot (were opt-in/OFF through Rev~C), a real behavior
    change for this procedure's own \code{ARM} step on hardware with
    nothing wired to PF13; and \code{ARM}'s failure reporting was
    reworked to name the specific failed precondition (which channel,
    which interlock) instead of one generic error. Documented firing
    via the external trigger as an alternative to Step~8's manual
    \code{SHOT:STARt}. Updated
    Section~\ref{sec:troubleshooting} to match. \\
  E & 2026-09-22 & -- & Wire-protocol rename: the output/demand commands
    moved out of \code{PID:*} into \code{SOURce:*}/\code{SHOT:*}/
    \code{LOG:*}/\code{CHANnel:*} (clean cut, no aliases);
    \code{PID:*} now means only \code{PID:GAINS}/\code{PID:LOOPMODE}.
    All command names in this document updated accordingly. \\
  \hline
\end{longtable}

\section{Purpose}

This document specifies the procedure for building, flashing, and
verifying WHAM-XREX-PFMG474 firmware, and for confirming a board is
running the expected firmware version before it is put into service or
handed off.

\section{Scope}
\label{sec:scope}

Covers firmware build and flashing (Sections~\ref{sec:build}--\ref{sec:verify}),
the legacy PWM-table shot path (Section~\ref{sec:table-fire}, Rev~B), and,
as of Rev~C, the modern closed-loop shot-profile path
(Section~\ref{sec:profile-fire}) -- this firmware \textbf{drives real
HRTIM PWM output} the moment \code{FIRE} (legacy path) or
\code{SHOT:STARt} (modern path) is sent.

\textbf{Correction, Rev~C:} Rev~B stated this firmware had ``no ARM
precondition, no state machine, and no fault interlock.'' The ARM/
state-machine part was accurate for Rev~B's own scope (the legacy
\code{TABLE:*}/\code{FIRE} path \emph{still} has no \code{ARM} step and
no state machine, by design -- Section~\ref{sec:table-fire}'s caution
covers this) -- but the ``no fault interlock'' part was already
incomplete even at Rev~B (\code{FIRE} has always rejected a latched
fault with \code{ERR 6}, see Section~\ref{sec:table-fire}'s own
caution, now corrected). A full state machine
(\code{IDLE}/\code{ARMED}/\code{FIRING}/\code{FAULT}), an \code{ARM}
precondition, and a richer, per-category fault interlock set
(overcurrent, water/temp, Enerpro, enable-output) were added to the
\emph{modern} path between Rev~B and Rev~C -- see Section~\ref
{sec:profile-fire}. Treat any gate driver or power stage wired to these
pins as live before either procedure, per your lab's normal HV/gate-
driver procedure -- the modern path's interlocks protect against
several real fault conditions but this SOP still provides no substitute
for that normal procedure. This document still does not cover power-
stage bring-up or calibration.

\section{Prerequisites}

\subsection{Equipment}

\begin{center}
\begin{tabular}{p{1.5in} p{3.2in}}
  \hline
  \textbf{Item} & \textbf{Notes} \\
  \hline
  WHAM-XREX-PFMG474 board & Powered per its own hardware documentation
    (not covered here). \\
  ST-Link (or equivalent SWD debug probe) & Required for the bootstrap
    flash (Section~\ref{sec:bootstrap}) and for recovery if a board
    stops responding over serial entirely. \\
  USB-to-serial adapter & Wired to USART2: adapter TX~$\to$~PA3
    (USART2\_RX), adapter RX~$\to$~PA2 (USART2\_TX), GND~$\leftrightarrow$~GND. \\
  \hline
\end{tabular}
\end{center}

\noindent\textit{(TX\slash RX wiring above is crossed: adapter TX goes
to the MCU's RX pin and vice versa, as usual for a point-to-point serial
link.)}

\subsection{Software}

\begin{center}
\begin{tabular}{p{1.5in} p{3.2in}}
  \hline
  \textbf{Tool} & \textbf{Notes} \\
  \hline
  STM32CubeIDE & Project build and the bootstrap SWD flash. \\
  \code{stm32flash} & AN3155 USART bootloader client. macOS:
    \code{brew install stm32flash}. \\
  Python 3 + \code{pyserial} & Runs \code{python/wham\_serial\_flash.py}.
    \code{pip install pyserial}. \\
  \hline
\end{tabular}
\end{center}

\caution{This firmware has not yet been through an ESD-controlled
production process. Use standard ESD precautions (grounding strap,
static-safe surface) when handling the bare board.}

\section{Procedure: Build the Firmware}
\label{sec:build}

\begin{enumerate}[label=\arabic*.]
  \item Open the project in STM32CubeIDE and press \textbf{Refresh}
    (\code{F5}) on the project if any files were added outside the IDE
    since it was last opened.
  \item Build the \textbf{Debug} configuration (\code{Project~$\to$~Build
    Project}, or from the command line: \code{cd Debug \&\& make all -j4}).
  \item Confirm the build produced no warnings. A clean build produces
    \code{WHAM-XREX-PFMG474.elf} in \code{Debug/}.
  \item \textbf{Generate the \code{.bin}.} This project's build
    configuration does not have the ``Convert to binary file'' post-build
    step enabled, so it must be generated by hand:
\begin{quote}\ttfamily\small\raggedright
arm-none-eabi-objcopy -O binary \textbackslash\\
\hspace*{1em}Debug/WHAM-XREX-PFMG474.elf \textbackslash\\
\hspace*{1em}Debug/WHAM-XREX-PFMG474.bin
\end{quote}
\end{enumerate}

\section{Procedure: Bootstrap Flash (SWD)}
\label{sec:bootstrap}

Required the \emph{first} time a board receives firmware with the
serial-bootloader feature (\code{BOOT} command), since that feature is
what the serial reflash path in Section~\ref{sec:serial-flash} depends
on -- a board cannot use it to receive the very first firmware that adds
it. Also use this procedure to recover a board that no longer responds
to \code{BOOT} over serial at all.

\begin{enumerate}[label=\arabic*.]
  \item Connect the ST-Link to the board's SWD header and to the host
    computer.
  \item In STM32CubeIDE, press \textbf{Run} (or \textbf{Debug}) to build
    and flash via the ST-Link, as normal for this IDE.
  \item Proceed to Section~\ref{sec:verify} to confirm the board is
    running the expected firmware.
\end{enumerate}

\section{Procedure: Serial Reflash}
\label{sec:serial-flash}

Use this for all firmware updates \emph{after} the bootstrap flash in
Section~\ref{sec:bootstrap} has been done once. This is the
remote-friendly path: no physical BOOT0\slash NRST\slash SWD access is
required, only the USART2 serial link.

\begin{enumerate}[label=\arabic*.]
  \item With the board powered and running its current firmware, and the
    USB-serial adapter connected, run:
\begin{quote}\ttfamily\small\raggedright
python3 python/wham\_serial\_flash.py \textbackslash\\
\hspace*{1em}--port /dev/cu.usbserial-XXXXX
\end{quote}
    substituting the actual serial device path. (On macOS, a single
    USB-serial adapter shows up as both a \code{cu.*} and \code{tty.*}
    device node; the script sees this as two candidates and will not
    auto-detect -- \code{--port} must be passed explicitly. Prefer the
    \code{cu.*} node.)
  \item The script will: send \code{BOOT} and wait for the
    acknowledgement; run \code{stm32flash} to write and verify
    \code{Debug/WHAM-XREX-PFMG474.bin}; then start the new firmware
    automatically (no manual reset needed -- see
    Section~\ref{sec:reset-note}).
  \item Confirm the script reports \texttt{[done]} with firmware
    written and verified before proceeding.
  \item Proceed to Section~\ref{sec:verify}.
\end{enumerate}

\subsection{Note on automatic restart after flashing}
\label{sec:reset-note}

Early testing of this procedure found that the newly-flashed firmware
could come up unresponsive over serial until the board was manually
power-cycled, because the STM32 ROM bootloader's \emph{Go} command (used
to start the new firmware without a physical reset) does not perform a
true chip reset, and left the processor's vector table pointer aimed at
the bootloader's own vector table rather than the application's. This
has been fixed in firmware (\code{boot\_jump.c}) -- a manual reset should
not be necessary. If a board flashed under this SOP does not respond to
the verification step below without a manual reset, that is a
regression and should be reported, not treated as expected behavior.

\section{Procedure: Verification}
\label{sec:verify}

\begin{enumerate}[label=\arabic*.]
  \item Open a serial connection to the board at \textbf{115200 8N1}
    (raised from 9600 as of Rev~B -- see \code{docs/changelog.txt}).
  \item Send \code{*IDN?} followed by a line ending (\code{\textbackslash
    r}, \code{\textbackslash n}, or both).
  \item Confirm a reply of the form:
    \begin{quote}
      \code{OK <board name> <hardware rev> <firmware version>}
    \end{quote}
    e.g.\ \code{OK WHAM-XREX-PFMG474 REVA v0.1}.
  \item Confirm the reported firmware version matches the version that
    was just built\slash flashed.
\end{enumerate}

Acceptance criterion for this SOP: Step 3 above returns a well-formed
reply with the expected version, on the first query, with no manual
reset performed since the flash in Section~\ref{sec:build} or
\ref{sec:serial-flash}.

\section{Procedure: Upload and Fire a PFM Table}
\label{sec:table-fire}

New in Rev~B. Uploads a complete PWM shot profile (a sequence of
period/duty entries, one per switching cycle) and plays it back on
demand. Full protocol detail: \code{docs/command\_reference.md}'s
\code{TABle:*} and \code{FIRE} entries.

\caution{\code{FIRE} drives real HRTIM PWM output on six pins (three
complementary phase pairs, $\sim$100\,ns hardware dead time between
each pair) the instant it is sent. There is still no \code{ARM} step
and no state machine on \emph{this legacy path} as of Rev~C -- \code{FIRE}
always takes effect immediately once its two checks pass (below),
whether the board was idle or a previous shot was still running; the
\code{ARM}/state-machine system added since Rev~B applies only to the
modern \code{SHOT:*} path (Section~\ref{sec:profile-fire}), not
to \code{FIRE}. \textbf{Correction, Rev~C:} \code{FIRE} DOES check for a
latched fault (\code{ERR 6} if one is, \code{FAULT:CLEAR} first) in
addition to the empty-table check below -- Rev~B's ``no fault
interlock'' claim for this path was incomplete even at the time and is
corrected here; this is still a single blanket check (any latched
fault, not per-category), not the modern path's richer interlock set.
Treat any gate driver or power stage wired to these pins as live before
proceeding, per your lab's normal HV/gate-driver procedure.}

\begin{enumerate}[label=\arabic*.]
  \item Complete Sections~\ref{sec:build}--\ref{sec:verify} first --
    a board must already be running current, verified firmware.
  \item Table \emph{construction} (what shot profile to build) lives
    entirely on the host, not in firmware -- see \code{pfm.h}'s own
    header note for why. \code{python/pfm\_table\_upload.py} is the
    reference builder: edit its \code{PROFILE} constant (1--5; see the
    script's own header comment for what each one is) before running:
\begin{quote}\ttfamily\small\raggedright
python3 python/pfm\_table\_upload.py \textbackslash\\
\hspace*{1em}--port /dev/cu.usbserial-XXXXX
\end{quote}
  \item Confirm the script reports \code{[done] table uploaded and
    confirmed} with the expected entry count (\code{TABLE:END -> OK
    <count>}).
  \item Optionally confirm independently at any time with \code{TABLE?}
    over a serial terminal (e.g.\ \code{python/scpi.py}) -- replies
    \code{OK <count>}, matching the upload.
  \item With the CAUTION above acknowledged and anything downstream of
    these six pins confirmed safe to drive, send \code{FIRE}. A prompt
    \code{OK} confirms the shot started; output stops on its own, at a
    coherent period boundary, once the table is exhausted -- there is
    no separate \code{STOP} command and no completion notification.
  \item Sending \code{FIRE} against an empty table (nothing uploaded
    yet) is rejected with \code{ERR 5} and drives nothing; sending it
    with a fault latched is rejected with \code{ERR 6} (send
    \code{FAULT:CLEAR} first) -- these two checks are this path's only
    built-in safeguards, not a substitute for the interlock discussed
    above.
\end{enumerate}

Acceptance criterion for this SOP: \code{TABLE:END} reports the exact
entry count the host script intended to send, and \code{FIRE} returns
\code{OK} with no \code{ERR}. Confirming the actual output waveform
(dead time, phase spacing, duty accuracy) is a bench/instrumentation
task outside this SOP's scope -- see \code{docs/test\_protocol.md}'s T4
group if that level of verification is needed.

\section{Procedure: Program and Fire a Shot Profile (Raw SCPI)}
\label{sec:profile-fire}

New in Rev~C. The modern shot path -- \code{ARM}/\code{SHOT:*} --
driven entirely by hand over a bare serial terminal, for an operator who
wants to understand (or reproduce without any host tooling) exactly what
\code{python/wham\_console.py}'s guided \code{shot} wizard does under the
hood. \textbf{For routine bench work, use the \code{shot} wizard instead}
(\code{python3 python/wham\_console.py}, then \code{shot} -- it applies
every step below automatically, with prompts, a pre-fire summary, and
automatic plotting) -- this procedure is for when that tooling is
unavailable, being debugged itself, or an operator specifically wants
direct command-by-command control. Full protocol detail for every
command below: \code{docs/command\_reference.md}'s \code{SOURce:*}/
\code{SHOT:*}/\code{LOG:*}/\code{CHANnel:*} and \code{XREX:*} entries.

\caution{\code{SHOT:STARt} (Step~8a) drives real HRTIM PWM output
on every enabled channel, ramping toward that channel's own commanded
demand current, from the moment it returns \code{OK}. Unlike the legacy
path (Section~\ref{sec:table-fire}), this path DOES have a state machine
(\code{IDLE} $\to$ \code{ARMED} $\to$ \code{FIRING}, back to \code{IDLE}
automatically at shot end or on a fault) and DOES check real
preconditions before \code{ARM} succeeds (Step~7) -- but those
preconditions are about channel readiness and interlock state, not
about whether a human has actually confirmed the downstream hardware is
safe to drive. Treat any gate driver or power stage wired to an enabled
channel's output as live before Step~8, per your lab's normal
HV/gate-driver procedure, exactly as for the legacy path.

\textbf{If firing via the external trigger instead (Step~8b, Rev~D):}
the board stays \code{ARMED} -- and will fire on the very next PF15
rising edge -- for as long as it takes an operator to reach for that
trigger, which may be well after Step~7. Treat \code{ARMED} with
\code{EXTernal:TRIGger} on as ``could fire at any moment,'' not just
``about to fire right now'': confirm the downstream hardware is safe
BEFORE sending \code{ARM} in this mode, not just before the trigger
pull, and \code{DISARM} promptly if the trigger will not be pulled
soon after arming.}

\begin{enumerate}[label=\arabic*.]
  \item Open a serial connection at \textbf{115200 8N1} and confirm the
    board responds to \code{*IDN?} (as in Section~\ref{sec:verify}).
  \item Confirm the starting state with \code{STATE?} -- expect
    \code{OK IDLE}. If it instead reports \code{OK FAULT <type>}, resolve
    the underlying condition first, then send \code{FAULT:CLEAR} (a
    fault cleared without resolving its real cause will simply re-latch,
    for every fault type this firmware detects).
  \item \textbf{For each channel to be used this shot}, send, in this
    order (every value explicit -- this firmware never resets a
    channel's configuration between shots on its own, so anything not
    sent here keeps whatever a previous shot, or a previous operator,
    left it at):
    \begin{enumerate}[label=(\alph*)]
      \item \code{XREX:CHANnel:ENAOut <ch> 1}
      \item \code{XREX:CHANnel:CONTactOut <ch> 1}
      \item \code{SOURce:ENAble <ch> 1}
      \item \code{PID:LOOPMODE <ch> <0|1>} -- \code{0} = open-loop
        (profile value written straight to output, no correction),
        \code{1} = closed-loop.
      \item \emph{Closed-loop only:} \code{PID:GAINS <ch> <Kp> <Ki> <Kd>}.
      \item \code{SHOT:CURRent <ch> <demandCurrentA>} -- this
        channel's peak current for the shot (clamped to its own
        configured full-scale current).
    \end{enumerate}
  \item \textbf{For every channel NOT being used this shot}, explicitly
    disable it: \code{SOURce:ENAble <ch> 0}, \code{XREX:CHANnel:
    ENAOut <ch> 0}, \code{XREX:CHANnel:CONTactOut <ch> 0}. Step~8 drives
    \emph{every currently-enabled} channel, not just the ones just
    configured in Step~3 -- a channel left enabled from an earlier shot
    will fire too, toward whatever demand current it was last given.
  \item Set the shared shot timing (one synchronized clock for every
    channel): \code{SHOT:TIMing <rampUpTimeS> <flatTopTimeS>
    <rampDownTimeS>} -- three independent durations as of Rev~C (all
    three required, each must be $>0$; ramp up and ramp down no longer
    have to match).
  \item \emph{Optional but recommended:} arm waveform logging, so the
    shot's actual behavior can be inspected afterward:
    \code{LOG:ARM <ch> <maxSamples> <decim>} for one channel, or
    \code{LOG:ARM 0 <maxSamples> <decim>} for every enabled channel at
    once, from the same real ticks.
  \item Send \code{ARM}. A plain \code{OK} confirms the transition to
    \code{ARMED}. On failure, the reply now names the specific
    precondition that failed (reworked Rev~D -- previously one generic
    \code{ERR 13} covered every cause): \code{ERR 13 Can't ARM -- not
    currently IDLE} (clear a fault or \code{DISARM} first);
    \code{ERR 15 External enable interlock not satisfied -- PF13 reads
    LOW} (see the \code{EXTernal:ENAble} note below); or
    \code{ERR 13 Channel~N's ENA\_OUT/CONTACT\_OUT not both set} (fix
    that channel per Step~3a/3b, or disable it per Step~4). Only the
    \emph{first} failing precondition is reported per attempt -- fix it
    and re-send \code{ARM} to see the next one, if any. Send
    \code{DISARM} at any point before Step~8 to cancel back to
    \code{IDLE} without firing anything.

    \textbf{\code{EXTernal:ENAble} default, Rev~D:} this interlock
    (PF13) now defaults \textbf{ON} at every boot (was opt-in/OFF
    through Rev~C). On hardware with a real external-enable signal
    wired to PF13, this Just Works. On a bare bench setup with nothing
    wired there, PF13's resting level depends on the specific hardware
    -- some fiber receivers idle HIGH with no signal present (interlock
    satisfied by default), a genuinely floating pin reads LOW via the
    board's own pull-down (interlock blocks \code{ARM} until either the
    real signal is wired, looped back via \code{DIAGnostic:GPOut11}
    (below), or \code{EXTernal:ENAble 0} is sent first for this
    session). Check \code{EXTernal:INPut?} to see PF13's actual current
    level before assuming which case applies.
  \item With the CAUTION above acknowledged, fire the shot -- either
    manually (send \code{SHOT:STARt}) or, \textbf{new Rev~D}, by
    physical trigger:
    \begin{enumerate}[label=(\alph*)]
      \item \emph{Manual:} send \code{SHOT:STARt}. \code{OK}
        confirms the shot has begun -- every enabled channel begins its
        own linear ramp (0\,A $\to$ its own \code{demandCurrentA} $\to$
        flat-top $\to$ 0\,A) from the same synchronized instant.
        \code{ERR 12} means \code{SHOT:TIMing} was never
        successfully sent this boot (Step~5); \code{ERR 15} means the
        external-enable interlock dropped between \code{ARM} and this
        command.
      \item \emph{External trigger:} send \code{EXTernal:TRIGger 1}
        (defaults ON as of Rev~D, same as \code{EXTernal:ENAble} above
        -- confirm with \code{EXTernal:TRIGger?} if unsure). While
        \code{ARMED}, a genuine \textbf{LOW-to-HIGH transition} on PF15
        fires the shot exactly as \code{SHOT:STARt} would --
        \textbf{edge-triggered, not level-triggered}: a signal already
        HIGH at the moment \code{ARM} succeeded does \emph{not} count,
        even if it stays HIGH afterward; only a fresh transition after
        arming fires it. Check \code{EXTernal:TRIGger:INPut?} to see
        PF15's current level, and see Section~\ref{sec:telemetry} for
        how to confirm the fire actually happened (a triggered shot
        emits its own \code{!EVT ... TRIGGER FIRING} line, distinct
        from a command-initiated one) -- unlike Step~(a), a trigger that
        fails to fire (wrong state, timing not set, interlock dropped)
        reports nothing back over serial; it just silently stays
        \code{ARMED}.
    \end{enumerate}
  \item Wait for completion: poll \code{SOURce:STATus? <ch>} (its
    \code{running} field drops to \code{0}), or simply wait
    \code{rampUpTimeS + flatTopTimeS + rampDownTimeS} plus a small
    margin. The shot ends itself (a full stop, output off entirely --
    not a hold-at-floor) with no separate \code{STOP} command and no
    completion notification; the state machine returns to \code{IDLE}
    on its own.
  \item If logging was armed, fetch it: \code{LOG:DATA? <ch>} --
    replies \code{OK <count> <rateHz>} then \code{count} groups of
    \code{setpointHz measuredHz outputHz}, one per logged sample.
  \item Return the board to a known-safe state before the next shot --
    Step~9's automatic stop does \emph{not} do this: for every channel
    used, \code{SOURce:ENAble <ch> 0}, \code{XREX:CHANnel:ENAOut
    <ch> 0}, \code{XREX:CHANnel:CONTactOut <ch> 0}, \code{PID:LOOPMODE
    <ch> 0}, \code{PID:GAINS <ch> 0 0 0}.
\end{enumerate}

Acceptance criterion for this SOP: \code{SHOT:STARt} returns
\code{OK} with no \code{ERR}; \code{SOURce:STATus?} shows \code{running=1}
immediately afterward and \code{running=0} once the total shot duration
has elapsed; and, if logging was armed, \code{LOG:DATA?} reports the
same sample count that was requested (or fewer only if the shot itself
was shorter than the log's own capacity). Confirming the actual output
waveform against the commanded profile (ramp linearity, flat-top
accuracy, closed-loop convergence) is a bench/instrumentation task
outside this SOP's scope -- \code{python/wham\_console.py}'s own
\code{plot}/\code{report} commands, or \code{python/run\_simulator\_
validation.py} against the bench simulator, cover that.

\section{Procedure: Telemetry -- Reading and Translating the Event Stream}
\label{sec:telemetry}

New in Rev~D. The controller emits \textbf{unsolicited} \code{!EVT}
lines -- device-initiated, not sent in response to any command -- any
time the top-level state (\code{IDLE}/\code{ARMED}/\code{FIRING}/
\code{FAULT}) changes, so an operator (or a host script) can observe
what the board is doing in real time instead of having to continuously
poll \code{STATE?}/\code{FAULT?}. A separate, always-on retained log
(the \textbf{flight recorder}) keeps the same history for post-mortem
review after the fact, even if nobody was watching the live stream when
it happened. Full protocol detail: \code{docs/command\_reference.md}'s
``Unsolicited telemetry events'' section and \code{docs/telemetry.md}.

\subsection{The live stream}

Enabled by default (\code{SYS:EVENT?} reports \code{1}). A \code{!EVT}
line can arrive interleaved between any two commands' own \code{OK}/
\code{ERR} replies -- the \code{!EVT} prefix can never collide with
either, so it is always unambiguous which kind of line just arrived.
\code{python/wham\_console.py} handles this automatically: every
\code{!EVT} line is printed inline, tagged with the host's own
wall-clock time, the moment it arrives -- nothing special has to be
sent or typed to see it happen. Driving the board over a bare terminal
instead, the same lines simply appear on their own between typed
commands and replies; no special terminal mode is needed. To silence
the live stream without losing anything (the flight recorder, below,
keeps recording regardless): \code{SYS:EVENT 0}. Re-enable with
\code{SYS:EVENT 1}.

\subsection{Translating an \code{!EVT} line}

Every line has the same shape: \code{!EVT <tick\_ms> <rest>}.
\code{<tick\_ms>} is a monotonic millisecond counter since this boot
(the same clock \code{SYS:TIME?} reports on demand) -- \textbf{not} a
wall-clock time; this board has no RTC. To relate an event's timestamp
back to a real time of day, note the host's own wall-clock time
alongside a \code{SYS:TIME?} query close in time to the event and work
from the difference -- \code{wham\_console.py}'s own inline tagging
(above) already does this automatically for the live stream.

\begin{center}
\begin{tabular}{p{2in} p{2.7in}}
  \hline
  \textbf{Line} & \textbf{Meaning} \\
  \hline
  \code{STATE IDLE} &
    Boot, a shot completed on its own, \code{DISARM}, \code{SOURce:STOP},
    or a successful \code{FAULT:CLEAR}. \\
  \code{STATE ARMED} &
    A successful \code{ARM} (Section~\ref{sec:profile-fire}, Step~7). \\
  \code{STATE FIRING} &
    Output has actually begun -- either \code{SHOT:STARt} sent
    manually, \emph{or} the external trigger (below). This line alone
    does \textbf{not} distinguish which. \\
  \code{TRIGGER FIRING} &
    Emitted as an EXTRA line, immediately alongside a \code{STATE
    FIRING} caused specifically by a PF15 rising edge
    (Section~\ref{sec:profile-fire}, Step~8b) -- \textbf{this} is how to
    tell a trigger-initiated shot apart from a command-initiated one: a
    \code{STATE FIRING} with no accompanying \code{TRIGGER FIRING} was
    started by \code{SHOT:STARt} (or \code{SOURce:RUN}); one
    with it was started by the physical trigger. \\
  \code{FAULT <type> [ch]} &
    A fault was just latched -- \code{<type>} is one of \code{GENERAL}
    (Water/Temp, or either hardware fault source, system-wide),
    \code{OVERCURRENT}, \code{ENABLE\_OUTPUT}, or \code{ENERPRO} (all
    three per-channel, \code{[ch]} 1-based and present only for these),
    or \code{EXTERNAL\_ENABLE} (system-wide, PF13 dropped while
    \code{FIRING}). See Section~\ref{sec:profile-fire}'s Scope
    correction and \code{docs/command\_reference.md} for what each type
    actually does to the output (ramp-down vs.\ immediate stop,
    system-wide vs.\ single-channel). \\
  \code{FAULT CLEAR} &
    A successful \code{FAULT:CLEAR} -- the underlying condition was
    confirmed actually gone, not merely requested; a \code{FAULT:CLEAR}
    that fails to actually clear (condition still present) produces no
    line at all, since nothing changed. \\
  \hline
\end{tabular}
\end{center}

\subsection{The flight recorder (post-mortem)}

\code{SYS:EVLOG?} replays every retained event since this boot, oldest
first, as a header \code{OK <n>} followed by \code{n} \code{!EVT}
lines in the exact same format as the live stream above -- the same
translation table applies unchanged. Retains up to 128 entries
(overwrite-oldest once full); RAM-only, so it clears on reset like
every other runtime state this firmware keeps. Critically, \textbf{this
keeps recording even with the live stream silenced} (\code{SYS:EVENT
0}) or with nobody connected/watching at all -- so the first
diagnostic step after ``something unexpected happened and I don't know
why'' should usually be \code{SYS:EVLOG?}, not a guess.

\textbf{Worked example.} A real session this same day: an operator
found the controller unexpectedly in \code{STATE?} \code{OK FAULT
GENERAL} with no memory of having caused it, and \code{GDS?} (the raw
gate-driver pin snapshot) showing every pin already healthy --
suggesting whatever tripped it had since cleared on its own, but the
\emph{latch} (by design, see Section~\ref{sec:table-fire}'s original
fault-interlock discussion) does not clear itself. \code{SYS:EVLOG?}
resolved it immediately:
\begin{verbatim}
> SYS:EVLOG?
< OK 3
< !EVT 40344 STATE ARMED
< !EVT 40354 STATE IDLE
< !EVT 275847 FAULT GENERAL
\end{verbatim}
This showed the fault fired spontaneously \textasciitilde{}4 minutes
\emph{after} the last ARM/IDLE activity, while genuinely \code{IDLE} --
ruling out a shot-in-progress cause and pointing instead at a real,
transient signal edge on the fiber-fed Water/Temp inputs (in this
case, traced to another board's own transmitter glitching during an
unrelated reset). \code{FAULT:CLEAR} then succeeded cleanly, since
\code{GDS?} had already confirmed the real condition was gone. Without
the flight recorder, the only available facts would have been ``it's
faulted now'' and ``the pins look fine now'' -- not enough, on their
own, to tell a real-but-resolved event apart from a false latch.

\section{Troubleshooting}
\label{sec:troubleshooting}

% longtable, not tabular/center: this table grew past one page as of
% Rev D (two new rows), and a plain tabular cannot break across a page
% boundary -- it either gets pushed whole onto the next page (leaving
% the previous one mostly blank) or overflows the bottom margin. \endhead
% repeats the header row on each continuation page.
\begin{longtable}{p{1.7in} p{3in}}
  \hline
  \textbf{Symptom} & \textbf{Cause / action} \\
  \hline
  \endhead
  \code{stm32flash} reports \code{Failed to init device, timeout.}
    before any write occurs &
    Not a firmware issue -- this probe doesn't depend on what's running
    on the chip. Check board power, serial wiring (TX\slash RX not
    crossed, common GND), and that no other program has the port
    open. \\
  \code{BOOT} gets no reply, but power/wiring check out &
    The firmware currently on the board may predate the \code{BOOT}
    command (run the bootstrap flash, Section~\ref{sec:bootstrap}) or
    may be a board affected by the \code{Go}-jump issue in
    Section~\ref{sec:reset-note} -- power-cycle once, then retry. \\
  Verification (Section~\ref{sec:verify}) fails after a serial reflash &
    Power-cycle the board once and repeat verification. If it then
    succeeds, this is the regression described in
    Section~\ref{sec:reset-note} -- report it. If it still fails, treat
    as a failed flash and repeat Section~\ref{sec:serial-flash} from the
    start. \\
  \code{ARM} returns \code{ERR 13}/\code{ERR 15}
    (Section~\ref{sec:profile-fire}) &
    \textbf{Reworked Rev~D: the reply itself now names the cause} --
    \code{ERR 13 Can't ARM -- not currently IDLE} (clear a fault or
    \code{DISARM} first); \code{ERR 15 External enable interlock not
    satisfied -- PF13 reads LOW} (see Step~7's own note, and
    \code{EXTernal:INPut?}); or \code{ERR 13 Channel~N's ENA\_OUT/
    CONTACT\_OUT not both set} (that channel's \code{XREX:CHANnel:
    ENAOut}/\code{CONTactOut} aren't both asserted --
    \code{SOURce:ENAble} alone is not sufficient). Only the first
    failing precondition is reported per attempt; fix it and re-send
    \code{ARM} to see the next one if any. \\
  \code{ARM} returns \code{ERR 15} unexpectedly, on hardware with
    nothing physically wired to PF13 &
    New Rev~D: \code{EXTernal:ENAble} now defaults \textbf{ON} at every
    boot. Whether this blocks \code{ARM} on unwired hardware depends on
    how that specific board's PF13 rests with nothing connected --
    check \code{EXTernal:INPut?} directly rather than assuming. If it
    reads LOW and no real interlock is available for this session, send
    \code{EXTernal:ENAble 0} first (see Step~7's own note). \\
  \code{SHOT:STARt} returns \code{ERR 12}
    (Section~\ref{sec:profile-fire}) &
    \code{SHOT:TIMing} was never successfully sent this boot, or
    was rejected -- as of Rev~C it takes three arguments (\code{SHOT:
    TIMing} itself will report the same \code{ERR 12} with
    ``needs three arguments'' if only two were sent, the pre-Rev-C
    form). \\
  A channel that should be silent shows real output or logged data
    (Section~\ref{sec:profile-fire}) &
    Confirm it was explicitly disabled THIS session (Step~4) --
    \code{SHOT:STARt} drives every \emph{currently-enabled}
    channel, not just the ones just configured, and enable state is
    never implicitly reset between shots or between operators. \\
  The external trigger (Step~8b) doesn't fire anything, with no error
    reported at all &
    New Rev~D. A triggered fire is a \emph{silent} no-op on failure --
    check, in order: (1) \code{STATE?} is \code{ARMED}, not some other
    state (only fires from \code{ARMED} -- see
    Section~\ref{sec:telemetry}'s translation table); (2)
    \code{SHOT:TIMing} was actually sent this boot (the same
    RAM-only config \code{SHOT:STARt} needs -- a fresh reflash
    wipes it, a real cause found live this same day); (3)
    \code{EXTernal:TRIGger:INPut?} is genuinely transitioning LOW then
    HIGH, not just sitting HIGH (edge-triggered, per Step~8b's own
    note); (4) the external-enable interlock, if on, is satisfied
    (\code{EXTernal:INPut?}). \code{SYS:EVLOG?}
    (Section~\ref{sec:telemetry}) settles which of these it was after
    the fact -- a genuine attempt that failed the interlock check still
    leaves no trace either way, exactly like the command-driven path's
    own re-check silently staying \code{ARMED}, so ruling out (1)-(3)
    first is the faster path. \\
  \hline
\end{longtable}

\section{References}

\begin{itemize}
  \item \code{AGENTS.md} -- project orientation and current functionality.
  \item \code{docs/command\_reference.md} -- full serial command protocol,
    including \code{TABle:*}/\code{FIRE} (Section~\ref{sec:table-fire})
    and \code{SOURce:*}/\code{SHOT:*}/\code{LOG:*}/\code{CHANnel:*}/
    \code{XREX:*} (Section~\ref{sec:profile-fire}).
  \item \code{docs/serial\_reflash\_guide.md} -- detailed reflash
    mechanism, including the \code{Go}-jump bug referenced in
    Section~\ref{sec:reset-note}.
  \item \code{docs/test\_protocol.md} -- bench-level verification,
    including waveform checks beyond this SOP's acceptance criteria.
  \item \code{python/pfm\_table\_upload.py} -- reference table builder
    used in Section~\ref{sec:table-fire}; see its own header comment
    for what each \code{PROFILE} value builds.
  \item \code{python/wham\_console.py} -- the maintained interactive
    operator console; its \code{shot}/\code{refire} commands apply
    every step of Section~\ref{sec:profile-fire} automatically, with
    guided prompts and automatic plotting -- the recommended tool for
    routine bench work, with Section~\ref{sec:profile-fire} as the
    reference for what it does under the hood.
  \item \code{docs/telemetry.md} -- the \code{!EVT} stream and flight
    recorder's own design writeup (wire format, buffering, the
    packet-framing invariant); Section~\ref{sec:telemetry} above is
    this SOP's operator-facing companion to it.
  \item \code{docs/changelog.txt} -- dated development history.
\end{itemize}

\end{document}
